\subsection{Quadratic Refinements over $\mathbb{F}_2$}

Throughout this subsection, vector spaces are finite-dimensional over $\mathbb{F}_2$.

\begin{definition} \label{def:arf-invariant-over-f2} \hlt{Quadratic Isometries and Arf Sums over $\mathbb{F}_2$}\\
Let $Q:\mathcal{V}\to\mathbb{F}_2$ be quadratic with nonsingular polar form $B_Q$. Let $\mathscr{B}_Q$ be the set of symplectic bases of $B_Q$.
\begin{enumerate}[label=(\roman*)]
\setlength{\itemsep}{0pt}
\item For quadratic $Q':\mathcal{W}\to\mathbb{F}_2$, \hlt{quadratic isometry} $Q\to Q'$ is linear isomorphism $\sigma:\mathcal{V}\to\mathcal{W}$ satisfying $Q'\circ\sigma=Q$. Write $Q\cong Q'$.
\item For $\beta=(\mathbf{u}_i,\mathbf{v}_i)_{i=1}^{m}\in\mathscr{B}_Q$, its \hlt{Arf sum} is $\operatorname{Arf}_{\beta}(Q)\defeq\sum\limits_{i=1}^{m}Q(\mathbf{u}_i)Q(\mathbf{v}_i)\in\mathbb{F}_2$.
\end{enumerate}
\end{definition}

\begin{example} \label{ex:characteristic-two-quadratic-refinements-finite-field} \hlt{One Polar Form with Two Refinements}\\
On $\mathcal{V}=\mathbb{F}_2^2$, set $H(x,y)=xy$ and $A(x,y)=x^2+xy+y^2$. Both have polar form $B((x,y),(u,v))=xv+yu$.\\
The identities $B((x,y),(0,1))=x$ and $B((x,y),(1,0))=y$ imply $\operatorname{rad}(B)=\{\mathbf{0}\}$.\\
Moreover, $H^{-1}(0)=\{(0,0),(1,0),(0,1)\}$, whereas $A^{-1}(0)=\{(0,0)\}$. A quadratic isometry preserves zero sets, so $H\not\cong A$.
\end{example}

\begin{theorem} \label{thm:f2-quadratic-refinements-arf-classification} \hlt{Arf Classification over $\mathbb{F}_2$}\\
Let $Q:\mathcal{V}\to\mathbb{F}_2$ be quadratic with nonsingular polar form $B_Q$.\\
For $m\in\mathbb{Z}_{\geq1}$ and $a\in\mathbb{F}_2$, define $Q_a^{(m)}(x_1,y_1,\ldots,x_m,y_m)\defeq x_1y_1+a(x_1^2+y_1^2)+\sum\limits_{i=2}^{m}x_iy_i$. Let $Q_0^{(0)}$ be the zero form on $\{\mathbf{0}\}$.
\begin{enumerate}[label=(\roman*)]
\setlength{\itemsep}{0pt}
\item $(\exists!\,m\in\mathbb{Z}_{\geq0})\,\dim\mathcal{V}=2m$.
\item For this $m$, $(\forall\,\beta\in\mathscr{B}_Q)\,Q\cong Q_{\operatorname{Arf}_{\beta}(Q)}^{(m)}$.
\item If $m\geq1$, then $(\forall\,a\in\mathbb{F}_2)\,\lvert(Q_a^{(m)})^{-1}(0)\rvert=2^{2m-1}+(-1)^a2^{m-1}$.
\item If $m\geq1$, then $Q_0^{(m)}\not\cong Q_1^{(m)}$.
\item $(\forall\,\beta,\gamma\in\mathscr{B}_Q)\,\operatorname{Arf}_{\beta}(Q)=\operatorname{Arf}_{\gamma}(Q)$; denote this common value by $\operatorname{Arf}(Q)$.
\item $Q\cong Q_{\operatorname{Arf}(Q)}^{(m)}$.
\item If $m\geq1$, then $(\forall\,a\in\mathbb{F}_2)\,\operatorname{Arf}(Q_a^{(m)})=a$.
\end{enumerate}
\begin{proof}
\begin{enumerate}[label=(\roman*)]
\setlength{\itemsep}{0pt}
\item The form $B_Q$ is alternating by Proposition \ref{prop:polar-forms-alternating-in-characteristic-two} and nonsingular by hypothesis.\\
By Theorem \ref{thm:symplectic-geometries-hyperbolic-classification}(i), $(\exists\,m\in\mathbb{Z}_{\geq0})\,\dim\mathcal{V}=2m$. Any such $m=\frac{1}{2}\dim\mathcal{V}$ is unique.
\item If $m=0$, then $Q=Q_0^{(0)}$ by $Q(\mathbf{0})=0$. Assume $m\geq1$ and fix $\beta=(\mathbf{u}_i,\mathbf{v}_i)_{i=1}^{m}\in\mathscr{B}_Q$. Put $r_i=Q(\mathbf{u}_i)$ and $s_i=Q(\mathbf{v}_i)$.\\
$(\forall\,x,y\in\mathbb{F}_2)\,Q(x\mathbf{u}_i+y\mathbf{v}_i)=r_ix^2+xy+s_iy^2$ by the polar identity and symplectic pairings.\\
The form $Q$ is the orthogonal sum of these plane restrictions by the vanishing cross-pairings.\\
If $(r_i,s_i)=(0,0)$, the $i$th summand is $H(x,y)\defeq xy$; if $(r_i,s_i)=(1,1)$, it is $A(x,y)\defeq x^2+xy+y^2$.\\
The forms $H$ and $A$ are quadratic with the standard symplectic polar form by direct calculation.\\
For $(r_i,s_i)=(0,1)$, the symplectic pair $(\mathbf{u}_i,\mathbf{u}_i+\mathbf{v}_i)$ has $Q$-values $(0,0)$.\\
For $(r_i,s_i)=(1,0)$, the symplectic pair $(\mathbf{u}_i+\mathbf{v}_i,\mathbf{v}_i)$ has $Q$-values $(0,0)$.\\
Every basis plane is $H$ or $A$ after these changes. The number of $A$-planes modulo $2$ is $\sum\limits_{i=1}^{m}r_is_i=\operatorname{Arf}_{\beta}(Q)$.\\
Fix two $A$-planes with symplectic bases $(\mathbf{e}_1,\mathbf{f}_1)$ and $(\mathbf{e}_2,\mathbf{f}_2)$. Set $\mathbf{p}_1=\mathbf{e}_1+\mathbf{e}_2$ and $\mathbf{q}_1=\mathbf{f}_1+\mathbf{e}_2$.\\
Set $\mathbf{p}_2=\mathbf{e}_1+\mathbf{f}_1+\mathbf{f}_2$ and $\mathbf{q}_2=\mathbf{e}_1+\mathbf{f}_1+\mathbf{e}_2+\mathbf{f}_2$.\\
The identity $(Q(\mathbf{p}_1),Q(\mathbf{q}_1),Q(\mathbf{p}_2),Q(\mathbf{q}_2))=(0,0,0,0)$ holds by direct calculation.\\
The equalities $B_Q(\mathbf{p}_1,\mathbf{q}_1)=B_Q(\mathbf{p}_2,\mathbf{q}_2)=1$ hold by direct calculation.\\
Moreover, $B_Q(\mathbf{p}_1,\mathbf{p}_2)=B_Q(\mathbf{p}_1,\mathbf{q}_2)=B_Q(\mathbf{q}_1,\mathbf{p}_2)=B_Q(\mathbf{q}_1,\mathbf{q}_2)=0$.\\
The tuple $(\mathbf{p}_1,\mathbf{q}_1,\mathbf{p}_2,\mathbf{q}_2)$ is a symplectic basis with zero $Q$-value on each basis vector by these identities.\\
Two $A$-planes are quadratically isometric to two $H$-planes by this basis change.\\
The isometry $Q\cong Q_{\operatorname{Arf}_{\beta}(Q)}^{(m)}$ holds by cancelling pairs of $A$-planes.
\item For $P:\mathcal{U}\to\mathbb{F}_2$, put $D(P)\defeq\sum\limits_{\mathbf{x}\in\mathcal{U}}(-1)^{P(\mathbf{x})}$.\\
If $P':\mathcal{U}'\to\mathbb{F}_2$ and $R(\mathbf{x},\mathbf{y})=P(\mathbf{x})+P'(\mathbf{y})$, then $D(R)=D(P)D(P')$ by factoring the two finite sums.\\
$D(H)=1+1+1-1=2$ and $D(A)=1-1-1-1=-2$ by direct calculation.\\
The form $Q_0^{(m)}$ is the orthogonal sum of $m$ copies of $H$ by definition.\\
The form $Q_1^{(m)}$ is the orthogonal sum of one $A$ and $m-1$ copies of $H$ by definition.\\
The equality $D(Q_a^{(m)})=(-1)^a2^m$ holds by multiplicativity, $D(H)=2$, and $D(A)=-2$. Put $N_a=\lvert(Q_a^{(m)})^{-1}(0)\rvert$.\\
By $D(Q_a^{(m)})=N_a-(4^m-N_a)=2N_a-4^m$, $N_a=2^{2m-1}+(-1)^a2^{m-1}$.
\item A quadratic isometry preserves zero-set cardinality by bijectivity.\\
The nonisometry $Q_0^{(m)}\not\cong Q_1^{(m)}$ holds by (iii), since the two cardinalities differ by $2^m>0$.
\item Fix $\beta,\gamma\in\mathscr{B}_Q$. If $m=0$, both Arf sums are empty and equal to zero.\\
For $m\geq1$, put $a=\operatorname{Arf}_{\beta}(Q)$ and $b=\operatorname{Arf}_{\gamma}(Q)$. The isometry $Q_a^{(m)}\cong Q_b^{(m)}$ holds by (ii).\\
The equality $a=b$ holds by (iv), because $a\neq b$ would identify this isometry with $Q_0^{(m)}\cong Q_1^{(m)}$.\\
The invariant $\operatorname{Arf}(Q)$ is well-defined by these two cases.
\item The isometry $Q\cong Q_{\operatorname{Arf}(Q)}^{(m)}$ holds by (ii) and (v).
\item Fix $a\in\mathbb{F}_2$. The map $Q_a^{(m)}$ is quadratic by direct calculation from its defining formula.\\
Its polar form has standard symplectic basis $\beta_a=(\mathbf{e}_i,\mathbf{f}_i)_{i=1}^{m}$ by the same calculation.\\
By definition, $Q_a^{(m)}(\mathbf{e}_1)=Q_a^{(m)}(\mathbf{f}_1)=a$. $(\forall\,i\in\mathbb{Z},\ 2\leq i\leq m)\,Q_a^{(m)}(\mathbf{e}_i)=Q_a^{(m)}(\mathbf{f}_i)=0$ by definition.\\
The equality $\operatorname{Arf}_{\beta_a}(Q_a^{(m)})=a^2=a$ holds by these values. By (v), $\operatorname{Arf}(Q_a^{(m)})=a$.
\end{enumerate}
\hlt{Structural Commentary}: Plane normalization and $A$-pair cancellation reduce $Q$; exponential sums separate the two residues.\\
\hlt{Key Insight}: The parity of anisotropic planes is the basis-independent quadratic residue.
\end{proof}
\end{theorem}
